\section{Budget-Constrained Pacing}
\label{sec:exp4_results}

Experiments~1--3 documented learning-driven bid suppression, in which reinforcement-learning and bandit agents converge to sub-competitive bids because they fail to reach competitive equilibria. This section studies a fundamentally different mechanism. Budget-constrained pacing agents suppress bids not because they fail to learn competitive strategies, but because they must ration a finite budget across auctions. We study this constraint-driven bid suppression using two pacing mechanisms, namely multiplicative dual pacing \citep{Balseiro2019Learning} in Experiment~4a and proportional-integral control in Experiment~4b. Both share a common environment. $N$ advertisers bid in $D = 100$ episodes of $T = 1{,}000$ auctions, with budgets resetting between episodes while control variables warm-start, mimicking daily budget cycles in ad exchanges. Valuations are drawn from bidder-specific log-normal distributions, introducing persistent asymmetry. Each experiment uses a $2^6$ full factorial replicated across 8 seeds (512 runs each). The theoretical worst-case Price of Anarchy under budget constraints is~2 for both auction formats \citep{Aggarwal2019, Gaitonde2023}.

\subsection{Multiplicative Dual Pacing (Experiment~4a)}
\label{sec:exp4a_results}

Experiment~4a deploys dual-based autobidding agents. The $2^6$ factorial crosses auction format (first-price vs.\ second-price), bidder objective (value-maximizer vs.\ utility-maximizer), market thickness ($N = 2$ vs.\ $N = 4$), budget multiplier ($m = 0.25$ tight vs.\ $m = 1.0$ loose), reserve price ($r = 0.0$ vs.\ $r = 0.3$), and value dispersion ($\sigma = 0.1$ low vs.\ $\sigma = 0.5$ high), replicated across independent seeds (\ExpFourANobs{} runs total).

\begin{figure}[H]
  \centering
  \includegraphics[width=\textwidth]{figures/e4a_trace}
  \caption{Representative learning trajectory for a single trial of Experiment~4a (first-price auction, value-maximizer objective, 2~bidders, 100 episodes $\times$ 1{,}000 rounds). Mean bid-to-value ratio per episode with the competitive benchmark (unconstrained Nash equilibrium $b/v = 0.5$).}
  \label{fig:e4a_trace}
\end{figure}

\subsubsection{Efficiency and Welfare}

Each episode produces a realised allocation and payments. Platform revenue is the total payments collected per episode. Liquid welfare follows \citet{Gaitonde2023}, defined as $W(x) = \sum_i \min(B_i, \sum_t x_{it} v_{it})$, capping each bidder's contribution at their budget. The effective Price of Anarchy divides the LP-relaxed offline optimum by realised liquid welfare. Values are at least~1, with values closer to~1 indicating higher efficiency.\footnote{The LP maximises $\sum_{t,i} v_{ti} x_{ti}$ subject to $\sum_i x_{ti} \le 1$ per round and $\sum_t v_{ti} x_{ti} \le B_i$ per bidder, with fractional allocations. A greedy integer allocation provides a tighter benchmark. We report the LP-based version as the more conservative test.} Tables~\ref{tab:exp4a_ranked_rev} and~\ref{tab:exp4a_ranked_poa} report the ranked significant effects.

\input{tables/exp4a_ranked_rev}

\begin{figure}[H]
  \centering
  \includegraphics[width=0.7\textwidth]{figures/e4a_main_rev}
  \caption{Experiment~4a: Main effects plot for platform revenue across the three factors.}
  \label{fig:e4a_main_rev}
\end{figure}

\input{tables/exp4a_ranked_poa}

\begin{figure}[H]
  \centering
  \includegraphics[width=0.7\textwidth]{figures/e4a_main_poa}
  \caption{Experiment~4a: Main effects plot for effective Price of Anarchy.}
  \label{fig:e4a_main_poa}
\end{figure}

All observed PoA values are well within the theoretical bound of~2 \citep{Aggarwal2019, Gaitonde2023}. The absence of a significant auction format effect on PoA contrasts with the revenue results. First-price bid suppression reduces payments to the auctioneer without proportionally distorting allocative outcomes. Allocative efficiency, the fraction of rounds in which the highest-value bidder wins, provides a complementary view. Under utility-maximizing objectives, agents achieve near-perfect efficiency. Under value-maximizing objectives with four bidders, efficiency drops substantially, with bid-to-value ratios exceeding 1.5.\footnote{Utility-maximizing efficiency ranges from \ExpFourAUtilEffLow{} to \ExpFourAUtilEffHigh{} across cells; value-maximizing four-bidder efficiency drops to \ExpFourAValFourEffLow--\ExpFourAValFourEffHigh, with bid-to-value ratios exceeding 1.5.} The objective$\times$n\_bidders interaction ($t = \ExpFourAAllocEffObjxNbidT$, $p \ExpFourAAllocEffObjxNbidPFmt$) captures this pattern, with the efficiency cost of value-maximizing objectives larger under four bidders than under two.

\subsubsection{Revenue and Budget Dynamics}

The budget multiplier dominates platform revenue ($|t| = \ExpFourARevFactorTOne$), shifting revenue by \ExpFourARevTopPctEffect\% of the grand mean. Tight budgets are associated with lower bids regardless of other design choices; loose budgets are associated with higher bids and greater competitive intensity. The number of bidders ranks fourth ($|t| = \ExpFourARevNbidAbsT$) and auction format fifth ($|t| = \ExpFourARevAuctionT$). Table~\ref{tab:exp4a_ranked_rev} reports the full hierarchy. The interaction plot (Figure~\ref{fig:e4a_int_rev}) reveals how factor combinations jointly shape revenue outcomes. Market thickness interacts with auction format, with the revenue gap between formats widening as the number of bidders increases.

\begin{figure}[H]
  \centering
  \includegraphics[width=0.7\textwidth]{figures/e4a_int_rev}
  \caption{Experiment~4a: Interaction plot for platform revenue.}
  \label{fig:e4a_int_rev}
\end{figure}

\subsubsection{Bidding Behaviour and Bid Suppression Indicators}

\input{tables/exp4a_ranked_vol}

\begin{figure}[H]
  \centering
  \includegraphics[width=0.7\textwidth]{figures/e4a_main_vol}
  \caption{Experiment~4a: Main effects plot for bid-to-value ratio.}
  \label{fig:e4a_main_vol}
\end{figure}

The factorial model explains \ExpFourAWelfareRsqPct\% of the variance in liquid welfare but only \ExpFourADriftRsqPct\% in cross-episode drift, reflecting substantial seed-level noise in the latter metric. Cross-episode drift, measured as the slope of the bid-to-value ratio across post-burn-in episodes, tests whether the bid-to-value ratio exhibits a downward trend across campaign cycles, as predicted by \citet{PaesLeme2024}. Although the factorial ANOVA for drift is statistically significant ($F$-test $p \ExpFourADriftFPFmt$), the model explains only \ExpFourADriftRsqPct\% of the variance ($R^2 = \ExpFourADriftRsq$), and all cell-level drift estimates remain near zero.\footnote{The statistical significance with low $R^2$ reflects the large sample size (\ExpFourANobs{} runs) detecting economically negligible effects.} This near-zero drift falsifies the prediction of \citet{PaesLeme2024} that dual-based pacing produces progressive bid suppression.

\subsubsection{Learning Dynamics}

Warm-start benefits are dominated by bidder objective ($t = \ExpFourAWarmObjT$). Under utility-maximizing objectives, warm-start produces measurable gains, while value-maximizing agents show negligible warm-start gains. Inter-episode volatility is similarly dominated by objective ($t = \ExpFourAIEVolObjT$), with utility-maximizing agents exhibiting higher volatility. Model adequacy diagnostics (Appendix~\ref{sec:appendix_robustness}) confirm clean inference.\footnote{$R^2$ ranges from \ExpFourARsqMin{} to \ExpFourARsqMax, PRESS gaps are below \ExpFourAPRESSGapMax, and \ExpFourABHPct\% of significant effects survive Benjamini--Hochberg correction.}

\subsubsection{Summary}

The budget multiplier dominates revenue (\ExpFourARevFactorOne{}, $|t| = \ExpFourARevFactorTOne$), with bidder objective ranking second and their interaction third. The number of bidders ranks fourth ($|t| = \ExpFourARevNbidAbsT$), and auction format is a comparatively minor factor ($|t| = \ExpFourARevAuctionT$). All PoA values fall well within the theoretical worst-case bound of~2. Global sensitivity analysis confirms budget multiplier dominance across all response variables (Appendix~\ref{sec:appendix_sensitivity}).

\subsection{PI Controller Pacing (Experiment~4b)}
\label{sec:exp4b_results}

Experiment~4b replaces the multiplicative dual pacing with a PI controller. The $2^6$ factorial crosses auction format, aggressiveness ($a_{\text{low}}$ vs.\ $a_{\text{high}}$), market thickness ($N = 2$ vs.\ $N = 4$), budget multiplier ($m = 0.25$ vs.\ $m = 1.0$), reserve price ($r = 0.0$ vs.\ $r = 0.3$), and value dispersion ($\sigma = 0.1$ vs.\ $\sigma = 0.5$), each replicated across 8 independent seeds (512 runs total).

\begin{figure}[H]
  \centering
  \includegraphics[width=\textwidth]{figures/e4b_trace}
  \caption{Representative learning trajectory for a single trial of Experiment~4b (first-price auction, default aggressiveness, 2~bidders, 100 episodes $\times$ 1{,}000 rounds). Mean bid-to-value ratio per episode with the competitive benchmark ($b/v = 1.0$).}
  \label{fig:e4b_trace}
\end{figure}

\subsubsection{Efficiency and Welfare}

The PI controller produces welfare outcomes that can be compared directly with the dual pacing results of Experiment~4a. Platform revenue and liquid welfare are computed identically, and the effective Price of Anarchy uses the same LP-based offline optimum as denominator. The budget multiplier again dominates revenue (Table~\ref{tab:exp4b_ranked_rev}), with the effect amounting to \ExpFourBRevTopPctEffect\% of the grand mean, confirming the pattern from Experiment~4a.

\input{tables/exp4b_ranked_rev}

\begin{figure}[H]
  \centering
  \includegraphics[width=0.7\textwidth]{figures/e4b_main_rev}
  \caption{Experiment~4b: Main effects plot for platform revenue.}
  \label{fig:e4b_main_rev}
\end{figure}

\input{tables/exp4b_ranked_poa}

\begin{figure}[H]
  \centering
  \includegraphics[width=0.7\textwidth]{figures/e4b_main_poa}
  \caption{Experiment~4b: Main effects plot for effective Price of Anarchy.}
  \label{fig:e4b_main_poa}
\end{figure}

For the effective Price of Anarchy, auction format is the dominant factor ($|t| = 10.7$), with first-price auctions producing higher PoA values. The number of bidders $\times$ budget multiplier interaction ranks second ($|t| = 9.9$). Thin markets with tight budgets generate the largest efficiency losses. Value dispersion reduces PoA ($|t| = 4.8$). All observed PoA values remain well within the theoretical worst-case bound of~2 \citep{Aggarwal2019, Gaitonde2023}, consistent with Experiment~4a. Allocative efficiency follows a similar pattern. The budget multiplier dominates ($|t| = 44.7$), with loose budgets yielding near-efficient allocations.

\subsubsection{Revenue and Budget Dynamics}

The auction type $\times$ budget multiplier interaction is the strongest two-way effect on revenue ($|t| = 9.5$), indicating that the revenue advantage of first-price auctions is concentrated in loose-budget settings where agents can bid aggressively. Under tight budgets, both formats converge as pacing constraints bind. The auction type $\times$ value dispersion interaction ($|t| = 5.6$) shows that the first-price revenue advantage is larger in high-$\sigma$ environments than in low-$\sigma$ environments. The budget multiplier $\times$ value dispersion ($|t| = 4.2$) and number of bidders $\times$ value dispersion ($|t| = 3.5$) interactions round out the significant two-way effects. Figure~\ref{fig:e4b_int_rev} displays these interaction patterns.

\begin{figure}[H]
  \centering
  \includegraphics[width=0.7\textwidth]{figures/e4b_int_rev}
  \caption{Experiment~4b: Interaction plot for platform revenue.}
  \label{fig:e4b_int_rev}
\end{figure}

\subsubsection{Bidding Behaviour}

The bid-to-value ratio is well-explained by the factorial model ($R^2 = 0.89$). The budget multiplier is the dominant factor ($|t| = 54.4$). Loose budgets permit higher bids relative to value, while tight budgets force aggressive shading. Auction format ranks second ($|t| = 22.3$), with first-price auctions producing substantially lower bid-to-value ratios than second-price, consistent with the theoretical prediction that first-price bidders shade below value. The auction type $\times$ budget multiplier interaction ($|t| = 17.3$) indicates that bid shading under first-price rules is amplified when budgets are tight. The auction type $\times$ number of bidders interaction ($|t| = 5.8$) reflects that additional bidders narrow the bid-to-value gap between formats. Compared with the multiplicative pacing of Experiment~4a, the PI controller produces qualitatively similar bid-to-value patterns but with different magnitudes. The PI controller produces less extreme shading under tight budgets than multiplicative pacing. Table~\ref{tab:exp4b_ranked_vol} and Figure~\ref{fig:e4b_main_vol} report the full effect hierarchy.

\input{tables/exp4b_ranked_vol}

\begin{figure}[H]
  \centering
  \includegraphics[width=0.7\textwidth]{figures/e4b_main_vol}
  \caption{Experiment~4b: Main effects plot for bid-to-value ratio.}
  \label{fig:e4b_main_vol}
\end{figure}

\subsection{Comparative Summary}

The comparison between multiplicative dual pacing and PI control isolates the role of the control mechanism from the budget constraint itself. In both experiments, the budget multiplier is the dominant revenue driver ($|t| = \ExpFourARevFactorTOne$ in Experiment~4a; $|t| = \ExpFourBRevFactorTOne$ in Experiment~4b), confirming that budget tightness, not the pacing algorithm, is the primary determinant of revenue outcomes. The number of bidders ranks among the top factors in both cases ($|t| = \ExpFourARevNbidAbsT$ in 4a; $|t| = 16.1$ in 4b), while auction format plays a stronger role under PI control ($|t| = 11.0$) than under dual pacing ($|t| = \ExpFourARevAuctionT$). All PoA values remain well within the theoretical worst-case bound of~2 across both pacing mechanisms. The two pacing mechanisms produce different bid-to-value ratio levels, but the factor hierarchy is consistent across both approaches, indicating that the structural features of the budget-constrained environment dominate the choice of control algorithm. This consistency further distinguishes constraint-driven from learning-driven bid suppression. In Experiments~1--3, switching algorithms changed the direction of the auction-format effect, whereas here, switching the pacing mechanism leaves the format effect unchanged. Global sensitivity analysis independently confirms these rankings. Budget multiplier achieves the highest total-order Sobol' index for revenue in both experiments ($S_T = 0.54$ in 4a, $S_T = 0.53$ in 4b). A notable difference emerges for the price of anarchy, where auction format is negligible under dual pacing ($S_T = 0.005$) but becomes the dominant factor under PI control ($S_T = 0.19$). Cross-method concordance is strong in both cases ($\rho = 0.72$ in 4a, $\rho = 0.77$ in 4b), indicating robust agreement across all seven sensitivity methods (Appendix~\ref{sec:appendix_sensitivity}).
